% macros.tex -- notation for the NM-ROM methods section.
% Symbol conventions (kept deliberately disjoint from the old submission, which
% overloaded K for both the discrete operator and a latent-adjacent quantity):
%   N  intervals per axis        n  interior unknowns
%   R  spatial bank width        k  latent dimension
%   M  number of retained weak tests    m  number of empirical-quadrature nodes

\newcommand{\Real}{\mathbb{R}}
\newcommand{\Nat}{\mathbb{N}}

% --- discrete objects -------------------------------------------------------
\newcommand{\Aop}{A}                    % fixed linear discrete operator
\newcommand{\bank}{G}                   % frozen spatial bank, n x R
\newcommand{\head}{h_\theta}            % nonlinear coefficient map, R^k -> R^R
\newcommand{\tests}{P}                  % weak test matrix, M x n
\newcommand{\rowscale}{\Lambda}         % diagonal row scaling
\newcommand{\lift}{u_{g}}               % boundary lift
\newcommand{\bcmask}{\mu}               % boundary mask / vanishing factor
\newcommand{\latent}{z}
\newcommand{\umanifold}{u}

% --- operators and norms ----------------------------------------------------
\newcommand{\jac}[1]{J_{#1}}
\newcommand{\dhead}{D\head}
\newcommand{\norm}[1]{\lVert #1 \rVert}
\newcommand{\normF}[1]{\lVert #1 \rVert_{F}}
\newcommand{\inner}[2]{\langle #1, #2 \rangle}
\newcommand{\T}{^{\top}}
\newcommand{\Tr}{\operatorname{tr}}
\newcommand{\diag}{\operatorname{diag}}
\newcommand{\rank}{\operatorname{rank}}
\newcommand{\argmin}{\operatorname*{arg\,min}}

% --- residuals --------------------------------------------------------------
\newcommand{\rfull}{r}                  % full discrete residual
\newcommand{\rweak}{r_{w}}              % weak (projected, row-scaled) residual
\newcommand{\station}[2]{\eta(#1,#2)}   % normalized-gradient stationarity measure

% --- placeholders for generated numbers -------------------------------------
% Every numeric value in this manuscript must come from a generator script that
% reads the run JSONs. \todo marks a slot that the generator has not yet filled.
\newcommand{\todo}[1]{{\color{red}\textbf{[\,#1\,]}}}
\newcommand{\gen}[1]{\todo{pending: #1}}

% Subscripted variants (the base macros already carry a subscript, so
% \head_a / \rweak_n would be a double subscript).
\newcommand{\headc}[1]{h_{\theta,#1}}   % one component of the head output
\newcommand{\rweakn}[1]{r_{w,#1}}       % the weak residual of time step #1
